% This document was generated by an LLM.

\documentclass{essay}

\title{Self-Check Questions:\\ Domain Changes in Algebraic Manipulation}
\author{}
\date{}

\begin{document}

\maketitle

Work each question before consulting the answer key at the end. Questions are numbered
continuously from~1 to~25; the key uses the same numbers. Where a question asks you to
solve an equation, also state \emph{where} any restriction comes from --- the original
equation's domain, or a step you introduced.

\section{Expressions versus functions}

\begin{enumerate}

\item In what precise sense are $\dfrac{(x+1)\,x}{x}$ and $x+1$ \emph{equal}, and in what
sense are they \emph{not}? Name the two settings in which the answers differ.

\item State the natural domain of $\dfrac{x^2-4}{x-2}$. Cancel the common factor; what is
the simplified expression and its natural domain? How do the two domains differ?

\item Define ``removable singularity.'' Write the limit that certifies the singularity of
$\dfrac{x^2-4}{x-2}$ at $x=2$ is removable, and give its value.

\item True or false, with justification: cancelling a common factor inside a single rational
expression can only ever \emph{shrink} its natural domain.

\item Two functions agree at every point of a common set but have different domains. Are they
the same function? Explain why the answer flips depending on whether you regard them as
functions or as elements of the field $\R(x)$.

\item Construct a rewriting of $y = x+2$ whose natural domain is $\R\setminus\{3\}$
--- that is, one that removes exactly the point at $x=3$.

\end{enumerate}

\section{Solving equations}

\begin{enumerate}[resume]

\item When you set out to solve an equation, what is the fixed invariant that every
manipulation must respect?

\item Distinguish a ``blind spot'' from an ``exclusion.'' Which of the two, if any, do you
carry into the final answer?

\item State, in one sentence, the diagnostic question to ask about any point that appears to
be excluded during a solution.

\item For each operation, say whether it can \emph{add} solutions, \emph{drop} solutions, or
neither: (a)~multiplying both sides by $D(x)$; (b)~dividing both sides by $D(x)$;
(c)~squaring both sides.

\item Explain why ``always keep the restriction'' is wrong, and why ``always discard the
restriction'' is also wrong. Give the type of example that defeats each rule.

\item When you write ``$\dfrac{x^2-1}{x-1} = x+1$ for $x \neq 1$,'' does the annotation
$x \neq 1$ qualify the left-hand side, the right-hand side, or the equality? Explain.

\item True or false, with justification: an extraneous root is always the symptom of an
arithmetic mistake.

\end{enumerate}

\section{Worked problems}

\begin{enumerate}[resume]

\item Solve $\dfrac{x}{x-3} = \dfrac{3}{x-3}$. Identify any extraneous candidate and trace
its restriction to a source.

\item Solve $x^2 = 4x$ two ways: by dividing both sides by $x$, and by factoring. Which
solution does the division lose, and why?

\item Solve $\sqrt{2x+3} = x$. Produce all candidates by squaring, check each in the
original equation, and name the extraneous one and the step that produced it.

\item You rewrite the equation $x - 2 = 5$ as $\dfrac{(x-2)(x+4)}{x+4} = 5$ and continue.
Does the rewrite change the solution? Which point does it flag, is that point excluded, and
what do you do with it?

\item Solve $\dfrac{1}{x-1} + \dfrac{1}{x+1} = \dfrac{2}{x^2-1}$. State the solution set and
explain it in terms of where the restriction originates.

\item For $\dfrac{x^2 - x}{x} = 1$, answer in order: (a)~natural domain of the left side;
(b)~its simplification; (c)~the candidate solution; (d)~whether any candidate is excluded;
(e)~the final solution set.

\end{enumerate}

\section{Diagnosing fallacies}

\begin{enumerate}[resume]

\item In the derivation $a=b \Rightarrow a^2=ab \Rightarrow a^2-b^2 = ab-b^2 \Rightarrow
(a+b)(a-b) = b(a-b) \Rightarrow a+b = b \Rightarrow 2=1$, identify the illegal step and the
exact value of the quantity being divided by.

\item Diagnose the fallacy
$-1 = \sqrt{-1}\cdot\sqrt{-1} = \sqrt{(-1)(-1)} = \sqrt{1} = 1$.
Which identity is misapplied, and what is its domain of validity?

\item A ``proof'' starts from $x = -3$, squares to get $x^2 = 9$, and concludes $x = 3$.
Explain the error using the taxonomy of risky steps.

\item Construct your own fallacious argument that $0 = 3$ by concealing a division by zero.
State plainly where the division by zero occurs.

\end{enumerate}

\section{Synthesis}

\begin{enumerate}[resume]

\item Explain why, in a typical calculus exercise, you can divide by expressions repeatedly
and almost never be burned. What conditions make a blind spot harmless?

\item Give one example of a domain restriction that is part of a \emph{function's identity},
and one example of a restriction that is merely an \emph{artifact of a solving step}. State
clearly which is which.

\end{enumerate}

\newpage

\section*{Answer key}

\begin{enumerate}

\item They are equal as \emph{formal rational expressions} --- the same element of the field
$\R(x)$, where cancelling a common factor is legal. They are \emph{not} equal as
\emph{functions}: the first has natural domain $\R\setminus\{0\}$, the second has
domain $\R$, and a function includes its domain.

\item Natural domain: $x \neq 2$. Cancelling, $\dfrac{(x-2)(x+2)}{x-2} = x+2$ for $x \neq 2$.
The simplified form $x+2$ has domain $\R$. The two differ only at $x=2$, where the
original has a hole that the simplified form fills.

\item A removable singularity is a point absent from the natural domain at which the
two-sided limit exists (and is finite), so the function can be patched to be continuous
there. Here $\displaystyle\lim_{x\to2}\frac{x^2-4}{x-2} = \lim_{x\to2}(x+2) = 4$.

\item False. Cancelling a common factor \emph{enlarges} the domain (it fills a hole). It is
\emph{introducing} a denominator that shrinks the domain.

\item As functions: no, they are not the same, because the domain is part of a function's
identity. As elements of $\R(x)$: yes, they are the same object, because that field
identifies fractions equivalent under cross-multiplication and never evaluates at a point.
Function equality demands equal domains and equal values; field equality demands only the
algebraic equivalence.

\item For instance $y = \dfrac{(x+2)(x-3)}{x-3}$. The factor $\dfrac{x-3}{x-3}$ equals $1$
except at $x=3$, where it is $0/0$; this removes exactly the point $x=3$.

\item The solution set of the \emph{original} equation. It is fixed before any manipulation;
each step is only a means of reading it.

\item A \emph{blind spot} is a point where a step is undefined or fails to preserve
equivalence, so the method yields no information there; you must test it in the original
equation. An \emph{exclusion} is a point genuinely outside the original equation's domain;
it cannot be a solution. You keep an exclusion (as forbidden); you never simply keep a blind
spot --- you check it.

\item ``Is this point excluded by the original equation itself, or only by a step I
introduced?''

\item (a)~Multiplying by $D(x)$ can \emph{add} solutions (at the zeros of $D$) but cannot
drop any. (b)~Dividing by $D(x)$ can \emph{drop} solutions (at the zeros of $D$) but cannot
add any. (c)~Squaring can \emph{add} solutions (it is not injective) but cannot drop any.

\item ``Always keep'' fails when a restriction is step-induced and the excluded point is in
fact a solution --- e.g.\ rewriting $x+1=1$ as $\dfrac{x(x+1)}{x}=1$ would wrongly report no
solution, though $x=0$ solves the original. ``Always discard'' fails when the restriction
belongs to the original equation's domain --- e.g.\ keeping $x=1$ as a solution of an
equation whose denominators are $x-1$ admits a value where the equation is undefined.

\item It qualifies the \emph{left-hand side} (and therefore the equality). The right side
$x+1$ is perfectly defined at $x=1$; the equation holds only where the left side is defined,
namely $x \neq 1$.

\item False. An extraneous root is the \emph{expected} byproduct of a step that does not
preserve equivalence (multiplying by a vanishing factor, squaring). The error would be
failing to check candidates and discard it, not its appearance.

\item Original domain: $x \neq 3$. Multiplying by $x-3$ gives $x = 3$, which is excluded by
the original domain. The restriction comes from the \emph{original} equation, so $x=3$ is a
genuine exclusion: it is extraneous and rejected. Solution set: empty.

\item Factoring: $x^2 - 4x = 0 \Rightarrow x(x-4) = 0 \Rightarrow x \in \{0, 4\}$. Dividing
by $x$ gives only $x = 4$, losing $x = 0$, because the division silently assumed $x \neq 0$
while $x=0$ was a genuine solution.

\item Squaring: $2x+3 = x^2 \Rightarrow x^2 - 2x - 3 = 0 \Rightarrow (x-3)(x+1) = 0
\Rightarrow x \in \{3, -1\}$. Check: $\sqrt{9} = 3$ holds; $\sqrt{1} = 1 \neq -1$ fails. The
root $x = -1$ is extraneous, introduced by the (non-injective) squaring step. Solution:
$x = 3$.

\item The rewrite does not change the solution: $x = 7$. It flags $x = -4$, a
\emph{step-induced} point, so it is a blind spot, not an exclusion. Checking it in the
original: $-4 - 2 = -6 \neq 5$, so $x=-4$ is not a solution anyway --- but the rule is that
had it satisfied the original, you would keep it. Solution set: $\{7\}$.

\item Original domain: $x \neq \pm 1$. The left side combines to
$\dfrac{(x+1)+(x-1)}{x^2-1} = \dfrac{2x}{x^2-1}$, so the equation becomes $2x = 2$, giving
$x = 1$. But $x = 1$ is excluded by the original domain, so it is extraneous. Solution set:
empty.

\item (a)~$x \neq 0$. (b)~$\dfrac{x(x-1)}{x} = x - 1$ for $x \neq 0$. (c)~$x - 1 = 1
\Rightarrow x = 2$. (d)~$x = 2$ is not excluded. (e)~Solution set: $\{2\}$.

\item The illegal step is dividing both sides by $a - b$. Because $a = b$, the quantity
$a - b$ is exactly $0$, so this is a division by zero.

\item The misapplied identity is $\sqrt{ab} = \sqrt{a}\,\sqrt{b}$, which is valid only for
$a, b \ge 0$ (nonnegative reals). With $a = b = -1$ it fails, so the second equality is
invalid.

\item Squaring is not injective: both $x = 3$ and $x = -3$ satisfy $x^2 = 9$, so squaring
\emph{added} the candidate $x = 3$. Concluding $x = 3$ keeps the extraneous root and
discards the true value $x = -3$; checking against the original $x = -3$ would have caught
it.

\item One construction: for any number $a$, both $a \cdot 0 = 0$ and $(a+3)\cdot 0 = 0$, so
$a \cdot 0 = (a+3)\cdot 0$. ``Cancelling the common factor $0$'' gives $a = a + 3$, hence
$0 = 3$. The division by zero is the cancellation of the factor $0$.

\item Because the expressions you divide by are generically nonzero on the domain of
interest, the values that would trigger a blind spot usually are not solutions, and the
extraneous-root check (when performed) passes trivially. A blind spot is harmless precisely
when the divisor is provably nonzero on the relevant domain, or when you check the flagged
candidates regardless.

\item Function-identity restriction: $f(x) = \dfrac{1}{x}$ has domain $x \neq 0$
permanently --- it is part of what $f$ \emph{is}. Solving-step artifact: rewriting the
equation $x + 1 = 1$ as $\dfrac{x(x+1)}{x} = 1$ introduces a temporary $x \neq 0$ that
belongs to the detour, not to the problem.

\end{enumerate}

\end{document}
